% !TeX root = ../tesis.tex

\subsection{Qué aporta \HaloTwo en este capítulo}
\label{subsec:stm-halo2-y-complejidad-que-aporta-halotwo-en-este-capitulo}

En este capítulo, \HaloTwo se utiliza como un sistema de aritmetización \textit{PLONKish}: transforma la relación abstracta $\STMRelation$ en un circuito organizado en filas, columnas y restricciones locales \cite{cryptoeprint:2019/953}. Las columnas contienen valores asignados por el \textit{prover}, constantes del circuito o entradas públicas; las restricciones locales verifican operaciones del circuito, mientras que los argumentos de permutación y los \textit{lookups} permiten expresar igualdades entre celdas y pertenencia a tablas. La prueba se construye mediante compromisos polinomiales, cuya instancia concreta utiliza \KZG.

La definición abstracta de \SNARK y la relación entre circuitos y restricciones algebraicas se presentan en la Definición~\ref{def:apendice-b-snark-snark-zksnark} y en las secciones iniciales del Apéndice~\ref{chap:aritmetizacion-y-construccion-de-un-snark}; el funcionamiento de \KZG se desarrolla en la Sección~\ref{sec:aritmetizacion-y-construccion-de-un-snark-kzg}. Esta sección se concentra únicamente en cómo la aritmetización concreta representa los cuatro chequeos de $\STMRelation$ y en las consecuencias estructurales de esa compilación para el circuito STM.

\subsection{Del protocolo \SNARK al circuito STM}
\label{subsec:stm-halo2-y-complejidad-del-protocolo-snark-al-circuito-stm}

Como se formaliza en la Definición~\ref{def:apendice-b-snark-snark-zksnark}, un \SNARK recibe una instancia pública $x$, un \textit{witness} $w$ y produce una prueba $\pi$ que el verificador puede comprobar sin reconstruir el cómputo completo. En este capítulo, la instancia pública está formada por los parámetros de la época, el mensaje y los compromisos publicados; el \textit{witness} contiene las evidencias de los firmantes, sus caminos de Merkle y los valores auxiliares necesarios para verificarlas.

El objetivo de esta sección no es volver a definir el protocolo \SNARK, sino describir cómo la relación $\STMRelation$ se traduce a las restricciones del circuito que luego procesa \HaloTwo.

\subsection{Representación del circuito en \HaloTwo}
\label{subsec:stm-halo2-y-complejidad-representacion-del-circuito-en-halotwo}

En \HaloTwo, un circuito se organiza como un \Layout: una disposición tabular de filas y columnas que permite reutilizar la misma estructura de restricciones a lo largo del dominio de evaluación. Cada fila representa una instancia local del cálculo, mientras que las columnas almacenan los valores que participan de esas restricciones.

Las columnas \Advice contienen valores asignados por el \textit{prover}; las columnas \Fixed contienen constantes determinadas por la configuración del circuito; y las columnas \Instance contienen las entradas públicas que recibe el verificador. Esta distinción es relevante para el circuito STM, porque las evidencias de los firmantes forman parte del \textit{witness}, mientras que el mensaje, los parámetros de época y los compromisos forman parte de la instancia pública.

\subsection{De la relación \texorpdfstring{\(\NP\)}{NP} a una aritmetización \textit{PLONKish}}
\label{subsec:stm-halo2-y-complejidad-de-la-relacion-np-a-una-aritmetizacion-plonkish}

Sea $\EvalDomain = \{\omega^0,\omega^1,\ldots,\omega^{n-1}\} \subseteq \Fq$ un subgrupo multiplicativo de orden $n$, y sea $\VanishingPoly(X) = \prod_{h \in \EvalDomain} (X - h) = X^n - 1 \in \poly{\Fq}$ su \textit{vanishing polynomial}. La aritmetización de \HaloTwo consiste en asociar a cada columna del \Layout un polinomio de grado menor que $n$ interpolado sobre \EvalDomain. En lo que sigue usamos estos nombres sólo para describir la aritmetización concreta, mientras que la relación abstracta continúa expresándose mediante $\PublicInputs$ y $\PrivateInputs$. Notamos:
\begin{itemize}[noitemsep]
	\item $\AdvicePoly{1}(X), \ldots, \AdvicePoly{a}(X)$ a los polinomios de columnas \Advice,
	\item $\FixedPoly{1}(X), \ldots, \FixedPoly{f}(X)$ a los polinomios de columnas \Fixed,
	\item $\InstancePoly{1}(X), \ldots, \InstancePoly{u}(X)$ a los polinomios de columnas \Instance.
\end{itemize}

Entonces, cada fila $0 \leq r < n$ impone identidades locales de la forma
$$G_j \left(\AdvicePoly{1}(\omega^r), \ldots, \AdvicePoly{a}(\omega^r), \FixedPoly{1}(\omega^r), \ldots, \FixedPoly{f}(\omega^r), \InstancePoly{1}(\omega^r), \ldots, \InstancePoly{u}(\omega^r)\right) = 0,$$
donde cada $G_j$ es un polinomio multivariado de bajo grado.

\subsection{Restricciones específicas de la relación STM}
\label{subsec:stm-halo2-y-complejidad-restricciones-especificas-de-la-relacion-stm}

Aplicado al circuito \STM, esto significa que los cuatro chequeos ya introducidos en $\STMRelation(\PublicInputs, \PrivateInputs) = 1$ deben reescribirse como igualdades locales sobre filas del \Layout. En particular:
\begin{itemize}[noitemsep]
    \item los chequeos de Merkle Path se traducen a iteraciones de $\Hash$ entre filas consecutivas;
    \item los chequeos de firma se traducen a restricciones sobre puntos de $\E(\Fq)$, escalares en $\Fr$ y desafíos recomputados;
    \item la lógica de lotería en $\Fr$ se representa con elementos de $\Fq$;
    \item la unicidad de índices se traduce a una cadena de desigualdades estrictas entre filas adyacentes.
\end{itemize}

\subsection{Compilación algebraica y verificación global}
\label{subsec:stm-halo2-y-complejidad-compilacion-algebraica-y-verificacion-global}

Todas estas familias de restricciones se agregan luego en una única identidad global. Más precisamente, si notamos \ConstraintNumerator a la combinación aleatorizada de \textit{gates}, condiciones de borde, consistencia con \textit{public inputs}, permutaciones y \textit{lookups}, entonces el \textit{prover} debe exhibir un polinomio $\QuotientPoly \in \poly{\Fq}$ tal que $\ConstraintNumerator(X) = \QuotientPoly(X) \cdot \VanishingPoly(X)$.

Se debe verificar que $\ConstraintNumerator(x) = 0$ para todo $x \in \EvalDomain$. Se utiliza \KZG como \PCS para los polinomios involucrados, de modo que esa identidad se pueda verificar mediante $\Open$ en pocos puntos aleatorios derivados con \FiatShamir.

\subsection{Argumentos de permutación, \textit{lookups} y columnas auxiliares}
\label{subsec:stm-halo2-y-complejidad-argumentos-de-permutacion-lookups-y-columnas-auxiliares}

\HaloTwo utiliza argumentos de permutación y de \textit{lookup} para expresar igualdades entre celdas y pertenencia a tablas, respectivamente. El tratamiento general de estos mecanismos pertenece a la descripción del protocolo y se remite al Apéndice B. En el circuito STM sólo nos interesa su consecuencia concreta: las permutaciones garantizan la consistencia entre valores reutilizados en distintas filas del \Layout, mientras que los \textit{lookups} permiten implementar los \textit{range checks} de la lotería.

En el circuito \STM, los \textit{lookups} se usan para \textit{range checks}. La comparación de lotería $e_t \le \LotteryTarget_{i_t}$ se interpreta sobre una representación en $\Fq$, y no sobre $\Fr$. Como $\Fq$ no tiene un orden total compatible con $\Fr$, el circuito compara representantes canónicamente elegidos en \([0,q)\). Cada elemento $x \in \Fq$ se descompone como $x = \LowPart{x} + 2^{127}\HighPart{x}$, con $0 \le \LowPart{x} < 2^{127}$, y $0 \le \HighPart{x} < 2^{128}$.

Estas cotas se fuerzan mediante \textit{lookups} sobre tablas de potencias de dos. Una vez validada la descomposición, la comparación entera se reduce a
$$x < y \iff \left(\HighPart{x} < \HighPart{y} \right) \;\vee\; \left(\HighPart{x} = \HighPart{y} \wedge \LowPart{x} < \LowPart{y}\right).$$

En consecuencia, el costo algebraico del circuito no queda determinado únicamente por las columnas explícitas: también depende de los recursos auxiliares que \HaloTwo introduce para verificar estas condiciones de consistencia global.

\subsection{Uso de crates de \texorpdfstring{\textnormal{\Midnight}}{Midnight}}
\label{subsec:stm-halo2-y-complejidad-uso-de-crates-de-midnight}

La descripción anterior corresponde a \HaloTwo/\KZG en abstracto. En \Mithril, sin embargo, la aritmetización concreta no se escribe directamente contra una API mínima de columnas y \textit{gates}, sino a través de las crates \texttt{midnight\_proofs} y \texttt{midnight\_zk\_stdlib}. Este detalle importa matemáticamente porque fija una transformación intermedia
$$\mathcal{C}_{\mathsf{Mid}}: \STMRelation \to \Circuit_{\mathsf{Halo2/KZG}},$$
que no es neutra: determina el \textit{backend}, la forma de los componentes especializados del circuito disponibles, la reutilización de columnas y el grado mínimo del dominio necesario para sintetizar el circuito.

En la implementación, la relación \STM se expresa como un objeto del tipo \texttt{Relation}, y luego se compila mediante $\texttt{MidnightCircuit::from\_relation(\&circuit)}$. Si denotamos por \MidnightSTMCircuit al circuito resultante, entonces su tamaño mínimo de dominio viene dado por un parámetro $\MinDomainK = \texttt{min\_k()}$, de modo que el dominio \EvalDomain efectivamente usado tiene cardinalidad $n = 2^{\MinDomainK}$.

Aún si formulamos $\STMRelation$ sin hacer referencia a columnas concretas, la complejidad real del sistema queda determinada por la compilación que hace \Midnight de esa relación a un \Layout \HaloTwo.

La particularidad más importante de esa compilación es que los distintos componentes no necesariamente aportan columnas disjuntas. Si denotamos por $\ArithCols, \EccCols, \PoseidonCols, \RangeCols$ a las necesidades de columnas de cada componente principal, la arquitectura de \Midnight no toma en general la suma
$$\ArithCols+\EccCols+\PoseidonCols+\RangeCols,$$
sino un \textit{pool} de columnas común cuyo tamaño $\AdviceCols$ denota la cantidad de columnas \Advice reservadas por la configuración, y cumple
$$\AdviceCols \approx \max\{\ArithCols, \EccCols, \PoseidonCols, \RangeCols \},$$
más algunos recursos adicionales de soporte. Esto justifica por qué no es correcto estimar el costo del \textit{prover} sumando ingenuamente columnas ``por componente''.

\subsection{Grado algebraico del circuito}
\label{subsec:stm-halo2-y-complejidad-grado-algebraico-del-circuito}

Diremos que un circuito de \HaloTwo tiene \emph{grado algebraico} a lo sumo $d$ si toda restricción/\textit{gate} local $G(X_1, \ldots, X_m) \in \Fq[X_1, \ldots, X_m]$ puede escribirse mediante un polinomio multivariado de grado total a lo sumo $d$. \MidnightSTMCircuit tiene grado algebraico $5$.

Esto es así porque la parte dominante del grado viene de \MidnightPoseidonHash. En efecto, la S-box de \MidnightPoseidonHash viene dada por $S(x) = x^5$. Por lo tanto, toda ronda que aplique esa S-box introduce naturalmente términos de quinto grado en las variables de estado del hash. Si denotamos por $\mathbf{s}^{(r)} = (s^{(r)}_1, \ldots, s^{(r)}_t)$ al estado de una ronda, la transición de una ronda completa adopta la forma
$$\mathbf{s}^{(r+1)} = \mathbf{M} \cdot \left((s^{(r)}_1)^5, \ldots, (s^{(r)}_t)^5 \right) + \mathbf{c}^{(r)},$$
donde $\mathbf{M}$ es una matriz $MDS$ y $\mathbf{c}^{(r)}$ es un vector de constantes de ronda.

En consecuencia, las restricciones que implementan esta transición pueden alcanzar grado local $5$. Por eso, cuando afirmamos que ``el circuito \MidnightSTMCircuit es de grado $5$'', nos referimos a la cota algebraica dominante de su sistema de restricciones, determinada en este caso por la S-box de \MidnightPoseidonHash.

La compilación concreta en \HaloTwo incorpora además restricciones auxiliares, como las asociadas a permutaciones y \textit{lookups}. Estas restricciones pertenecen a la representación del circuito y no forman parte de la definición de la S-box. En la práctica, la representación resultante puede requerir fragmentar \QuotientPoly en varios \textit{chunks} cuando su grado excede el rango soportado por la \SRS; esta es una consecuencia del \textit{backend} y de la configuración de la prueba, no una modificación de la S-box.

\subsection{Costo estructural del \textit{prover}}
\label{subsec:stm-halo2-y-complejidad-costo-estructural-del-prover}

Denotamos por \BaseCols al número de polinomios correspondientes a columnas visibles, y por $\AuxCols = \PermCols + \LookupCols + \QuotCols$ al número de polinomios auxiliares inducidos por permutación, \textit{lookups} y \textit{quotient}. Una cota heurística razonable para el costo del \textit{prover} es
$$\ProverCost = \BigO\left((\BaseCols + \AuxCols) \, n \log n\right) +  \BigO\left(\CommitCols \cdot \MSM(n)\right),$$
donde $\CommitCols$ es el número de compromisos \KZG efectivamente producidos. Esta fórmula no es un teorema, pero sí es un resumen razonable del costo: transformaciones tipo \FFT sobre dominios de tamaño $n$, construcción de polinomios auxiliares y \MSM grandes sobre la \SRS para \textit{batch commitments} en \KZG. La operación \MSM y su incidencia en el costo se describen en la Subsección~\ref{subsec:apendice-b-groth16-complejidad-de-groth}.

El circuito \MidnightSTMCircuit parece estar dominado por aritmética de curva elíptica sobre \Jubjub, así como evaluaciones repetidas de \PoseidonHash.
